\documentclass{ximera}

\input{../../../preamble.tex}


\definecolor{fvdnavy}{HTML}{071D43}
\definecolor{fvdcyan}{HTML}{20BFC6}
\definecolor{fvdcyanlight}{HTML}{EFFCFB}
\definecolor{fvdmagenta}{HTML}{F10D69}
\definecolor{fvdmagentalight}{HTML}{FFF1F7}
\definecolor{fvdtext}{HTML}{163044}





%=================================================
% Pedagogische omgevingen
% PDF: tcolorbox
% HTML: learning-box
%=================================================

\newenvironment{voorkennis}
{%
    \ifdefined\HCode
        \HCode{%
<section class="learning-box learning-box--prior">
<div class="learning-box__header">
<span class="learning-box__icon" aria-hidden="true"></span>
<span class="learning-box__title">Voorkennis</span>
</div>
<div class="learning-box__content">
        }%
        \begin{itemize}
    \else
       \begin{tcolorbox}[
    enhanced,
    breakable,
    colback=fvdcyanlight,
    colframe=fvdcyan,
    coltext=fvdtext,
    boxrule=0.5pt,
    leftrule=4pt,
    arc=4mm,
    outer arc=4mm,
    left=7mm,
    right=6mm,
    top=5mm,
    bottom=5mm,
    before skip=6mm,
    after skip=6mm,
    title=Voorkennis,
    coltitle=fvdcyan,
    fonttitle=\bfseries\Large,
    attach title to upper={\par\smallskip},
    boxed title style={
        empty,
        left=0mm,
        right=0mm,
        top=0mm,
        bottom=0mm
    },
    overlay={
        \draw[
            fvdcyan,
            line width=0.5pt
        ]
        ([xshift=42mm,yshift=-13mm]frame.north west)
        --
        ([xshift=-7mm,yshift=-13mm]frame.north east);
    }
]
        \begin{itemize}
    \fi
}
{%
    \end{itemize}
    \ifdefined\HCode
        \HCode{%
</div>
</section>
        }%
    \else
        \end{tcolorbox}
    \fi
}


\newenvironment{leerdoelen}
{%
    \ifdefined\HCode
        \HCode{%
<section class="learning-box learning-box--goals">
<div class="learning-box__header">
<span class="learning-box__icon" aria-hidden="true"></span>
<span class="learning-box__title">Leerdoelen</span>
</div>
<div class="learning-box__content">
        }%
        \begin{itemize}
    \else
\begin{tcolorbox}[
    enhanced,
    breakable,
    colback=fvdmagentalight,
    colframe=fvdmagenta,
    coltext=fvdtext,
    boxrule=0.5pt,
    leftrule=4pt,
    arc=4mm,
    outer arc=4mm,
    left=7mm,
    right=6mm,
    top=5mm,
    bottom=5mm,
    before skip=6mm,
    after skip=6mm,
    title=Leerdoelen,
    coltitle=fvdmagenta,
    fonttitle=\bfseries\Large,
    attach title to upper={\par\smallskip},
    boxed title style={
        empty,
        left=0mm,
        right=0mm,
        top=0mm,
        bottom=0mm
    },
    overlay={
        \draw[
            fvdmagenta,
            line width=0.5pt
        ]
        ([xshift=42mm,yshift=-13mm]frame.north west)
        --
        ([xshift=-7mm,yshift=-13mm]frame.north east);
    }
]
        \begin{itemize}
    \fi
}
{%
    \end{itemize}
    \ifdefined\HCode
        \HCode{%
</div>
</section>
        }%
    \else
        \end{tcolorbox}
    \fi
}


\addPrintStyle{../../..}

\title{Oefeningen --- Lineaire regressie}
\author{Ingmar Herreman}


\begin{document}

% FVD_CHAPTER_DATA_START
% Automatisch gegenereerd uit index.tex.
% Niet handmatig aanpassen.
\fvdchapterdata{3}{Verbanden onderzoeken met data}{9}
% FVD_CHAPTER_DATA_END






\maketitle



\begin{oefening}[Parameters interpreteren]{\oefB \ESS}
Voor een dataset werd het regressiemodel

\[
\hat y=-2{,}8x+74
\]

gevonden. $x$ wordt gemeten in ${}^{\circ}\mathrm C$ en $y$ in
$\mathrm{kWh}$.

\begin{question}
Wat is de eenheid van de richtingscoëfficiënt?

\begin{oplossing}
\[
\frac{\mathrm{kWh}}{{}^{\circ}\mathrm C}.
\]
\end{oplossing}
\end{question}

\begin{question}
Interpreteer $-2{,}8$.

\begin{oplossing}
Wanneer $x$ met $1\,{}^{\circ}\mathrm C$ toeneemt, daalt de door het
model voorspelde waarde van $y$ met $2{,}8\,\mathrm{kWh}$.
\end{oplossing}
\end{question}

\begin{question}
Wat betekent $74$ algebraïsch?

\begin{oplossing}
Voor $x=0\,{}^{\circ}\mathrm C$ voorspelt het model
$\hat y=74\,\mathrm{kWh}$.
\end{oplossing}
\end{question}
\end{oefening}


\begin{oefening}[Voorspellen met een model]{\oefB \ESS}
Voor een sensor geldt binnen het onderzochte bereik

\[
\hat U=0{,}42T+1{,}8,
\]

met $T$ in ${}^{\circ}\mathrm C$ en $U$ in volt.

\begin{question}
Voorspel $U$ bij $T=20\,{}^{\circ}\mathrm C$.

\begin{oplossing}
\[
\hat U=0{,}42\cdot20+1{,}8=10{,}2\,\mathrm V.
\]
\end{oplossing}
\end{question}

\begin{question}
Wat betekent het dakje in $\hat U$?

\begin{oplossing}
Het gaat om de door het regressiemodel voorspelde waarde, niet
noodzakelijk om de werkelijk gemeten waarde.
\end{oplossing}
\end{question}
\end{oefening}


\begin{oefening}[Interpolatie of extrapolatie?]{\oefB \ESS}
Een model is gebaseerd op metingen voor

\[
15\leq x\leq 80.
\]

Classificeer de volgende voorspellingen.

\[
x=20,\qquad x=57,\qquad x=80,\qquad x=95,\qquad x=4.
\]

\begin{oplossing}
$x=20$, $57$ en $80$ liggen binnen het gemeten bereik en zijn
interpolaties. $x=95$ en $x=4$ liggen buiten het bereik en zijn
extrapolaties.
\end{oplossing}
\end{oefening}


\begin{oefening}[Gemeten en voorspeld]{\oefB}
Voor een waarneming geldt

\[
y=53{,}7
\qquad\text{en}\qquad
\hat y=50{,}9.
\]

\begin{question}
Bereken het residu.

\begin{oplossing}
\[
e=y-\hat y=53{,}7-50{,}9=2{,}8.
\]
\end{oplossing}
\end{question}

\begin{question}
Ligt het meetpunt boven of onder de regressierechte?

\begin{oplossing}
Boven de regressierechte, want het residu is positief.
\end{oplossing}
\end{question}
\end{oefening}





\begin{oefening}[Groei van jongensbaby's]{\oefU \ESS}

Gebruik

\csv
{https://raw.githubusercontent.com/Ingmar-1986/wiskunde_cursus_2/main/modules/Data/Statistiek/lengtebabys.csv}
{lengtebabys.csv}

Open de dataset met geschikte ICT. Maak een spreidingsdiagram, bepaal
$r$ en voer een lineaire regressie uit. Je hoeft de gegevens niet
handmatig over te nemen.

\begin{question}
Beoordeel op basis van de puntenwolk of een lineair model binnen het
gemeten leeftijdsgebied redelijk lijkt.

\begin{oplossing}
De punten liggen vrij dicht rond een stijgende rechte band. Binnen
het beperkte gemeten leeftijdsgebied is een lineair model daarom een
redelijke eerste benadering.
\end{oplossing}
\end{question}

\begin{question}
ICT geeft ongeveer
\[
r=0{,}976
\]
en
\[
\hat L=1{,}877t+54{,}733.
\]
Interpreteer de richtingscoëfficiënt met eenheid.

\begin{oplossing}
Volgens het model neemt de voorspelde gemiddelde lengte binnen het
onderzochte gebied per extra levensmaand met ongeveer
\[
1{,}88\,\mathrm{cm}
\]
toe.
\end{oplossing}
\end{question}

\begin{question}
Voorspel de gemiddelde lengte op $10{,}5$ maanden en classificeer de
voorspelling.

\begin{oplossing}
\[
\hat L(10{,}5)
=1{,}876923\cdot10{,}5+54{,}7333
\approx74{,}44\,\mathrm{cm}.
\]
Dit is interpolatie omdat $10{,}5$ maanden binnen het gemeten interval
van $1$ tot $12$ maanden ligt.
\end{oplossing}
\end{question}

\begin{question}
Waarom is hetzelfde model ongeschikt om de gemiddelde lengte op
$20$ jaar te voorspellen?

\begin{oplossing}
Dat is een extreme extrapolatie. Menselijke groei blijft niet lineair
doorlopen tot volwassen leeftijd, zodat het model buiten het beperkte
meetgebied zijn betekenis verliest.
\end{oplossing}
\end{question}

\end{oefening}


\begin{oefening}[Hoogte en temperatuur in weerstations]{\oefU \ESS}

Gebruik

\csv
{https://raw.githubusercontent.com/Ingmar-1986/wiskunde_cursus_2/main/modules/Data/Statistiek/weerstations.csv}
{weerstations.csv}

Maak met ICT een spreidingsdiagram van temperatuur in functie van hoogte
en bepaal een lineair regressiemodel.

\begin{question}
Beschrijf de puntenwolk en beoordeel of een lineair model zinvol lijkt.

\begin{oplossing}
De puntenwolk vertoont een sterk negatief en ongeveer lineair verband:
hogere stations zijn in deze dataset doorgaans kouder.
\end{oplossing}
\end{question}

\begin{question}
ICT geeft ongeveer
\[
r=-0{,}975
\]
en
\[
\hat T=-0{,}00542H+10{,}995.
\]
Interpreteer de richtingscoëfficiënt liever per $100\,\mathrm m$.

\begin{oplossing}
Per extra $100\,\mathrm m$ hoogte daalt de door het model voorspelde
temperatuur gemiddeld met ongeveer
\[
0{,}542\,{}^\circ\mathrm C.
\]
\end{oplossing}
\end{question}

\begin{question}
Gebruik het model om de temperatuur op $1000\,\mathrm m$ te voorspellen.

\begin{oplossing}
\[
\hat T(1000)
=-0{,}0054232\cdot1000+10{,}9948
\approx5{,}57\,{}^\circ\mathrm C.
\]
Omdat $1000\,\mathrm m$ binnen het gemeten hoogtegebied ligt, is dit
interpolatie.
\end{oplossing}
\end{question}

\begin{question}
Waarom moet een voorspelling voor $8000\,\mathrm m$ kritisch worden
beoordeeld?

\begin{oplossing}
$8000\,\mathrm m$ ligt ver buiten het gemeten hoogtegebied. Dit is
een zeer verre extrapolatie en het lokale lineaire verband hoeft daar
niet geldig te blijven.
\end{oplossing}
\end{question}

\end{oefening}


\begin{oefening}[Archaeopteryx]{\oefU \ESS}

Gebruik

\csv
{https://raw.githubusercontent.com/Ingmar-1986/wiskunde_cursus_2/main/modules/Data/Statistiek/archaeopteryx.csv}
{archaeopteryx.csv}

De dataset bevat lengtes van dijbeen en opperarmbeen van fossielen.

\begin{question}
Maak een spreidingsdiagram en beoordeel de lineariteit.

\begin{oplossing}
De vijf punten liggen zeer dicht rond een stijgende rechte band.
Een lineair model lijkt voor deze beperkte dataset zeer geschikt.
\end{oplossing}
\end{question}

\begin{question}
Bepaal met ICT $r$ en de regressierechte voor opperarmbeen in functie
van dijbeen.

\begin{oplossing}
Ongeveer:
\[
r=0{,}994
\]
en
\[
\widehat{\text{opperarmbeen}}
=1{,}197\cdot\text{dijbeen}-3{,}660.
\]
\end{oplossing}
\end{question}

\begin{question}
Welke waarschuwing hoort bij de zeer hoge correlatie?

\begin{oplossing}
De dataset bevat slechts vijf waarnemingen. De hoge correlatie toont
een zeer sterke lineaire samenhang in deze gegevens, maar zegt niet
dat elk toekomstig fossiel exact op dezelfde rechte zal liggen.
\end{oplossing}
\end{question}

\end{oefening}




\begin{oefening}[Van puntenwolk naar modelkeuze]{\oefU \ESS}

Vergelijk

\begin{itemize}
  \item \texttt{lengtebabys.csv},
  \item \texttt{weerstations.csv},
  \item \texttt{druk.csv}.
\end{itemize}

Gebruik ICT om voor elke dataset een spreidingsdiagram en een
regressierechte te verkrijgen.

\begin{question}
Bij welke datasets lijkt een lineair regressiemodel grafisch het meest
verantwoord?

\begin{oplossing}
Bij \texttt{lengtebabys.csv} en \texttt{weerstations.csv} liggen de
punten ongeveer rond een rechte band. Bij \texttt{druk.csv} is het
verband sterk maar systematisch gekromd, zodat een rechte belangrijke
structuur mist.
\end{oplossing}
\end{question}

\begin{question}
Waarom mag software die altijd een regressierechte berekent niet zelf
de modelkeuze bepalen?

\begin{oplossing}
Software kan voor elke dataset een rechte berekenen. De leerling moet
echter op basis van puntenwolk, residuen en context beoordelen of die
rechte ook een zinvolle beschrijving is.
\end{oplossing}
\end{question}

\end{oefening}




\begin{oefening}[Wiskunde voorspellen uit fysica?]{\oefG}

Gebruik

\csv
{https://raw.githubusercontent.com/Ingmar-1986/wiskunde_cursus_2/main/modules/Data/Statistiek/wiskfys.csv}
{wiskfys.csv}

Bepaal met ICT eerst de regressie van fysica op wiskunde.

Een leerling draait daarna die formule algebraïsch om en gebruikt ze
om wiskunde uit fysica te voorspellen.

\begin{question}
Is dat automatisch dezelfde regressierechte als wanneer je vanaf het
begin wiskunde op fysica regresseert? Leg uit.

\begin{oplossing}
Nee. Bij lineaire regressie worden de verticale afwijkingen ten opzichte
van de gekozen afhankelijke variabele geminimaliseerd. Wanneer de rollen
van $x$ en $y$ worden omgewisseld, verandert dus ook het regressieprobleem.
De ene regressierechte algebraïsch omkeren is in het algemeen niet
hetzelfde als de andere regressie uitvoeren.
\end{oplossing}
\end{question}

\end{oefening}




\end{document}

